\section{Predictive Models and Graph-Based Learning}
\label{sec:graph_learning_finance}

Before considering graph-based models, it is important to distinguish the information supplied to a model from the flexibility with which that information is processed. Ridge regression provides a regularized linear benchmark that limits coefficient instability when predictors are correlated \parencite{hoerl1970ridge}. Gradient-boosted decision trees can instead learn nonlinear relationships and interactions without requiring them to be specified in advance, with XGBoost providing a scalable implementation of this approach \parencite{chen2016xgboost}. Comparing these model classes on identical feature sets therefore tests whether nonlinear flexibility improves prediction without attributing that improvement to the inclusion of additional ownership or network information.

Representing institutional ownership as a network, as motivated in Section~\ref{sec:ownership_networks}, introduces a separate modeling choice. Engineered network statistics can be added to conventional tabular models as additional predictors, whereas graph neural networks can learn representations directly from the graph. GraphSAGE constructs inductive node representations by aggregating information from neighboring nodes \parencite{hamilton2017inductive}. Graph attention networks (GATs) extend this approach by learning the relative importance assigned to different neighbors during aggregation \parencite{velickovic2018graph}. Applied to a manager--stock network, these methods allow a security's representation to depend on both its own characteristics and the managers connected to it, potentially capturing relational patterns that are not fully summarized by manually engineered network features.

Applications of graph learning to institutional holdings data are still limited and have emerged only recently. \textcite{izadifar2026institutional} represent Form 13F holdings as a discrete-time temporal bipartite graph connecting institutional managers and securities. They formulate the prediction of future portfolio weights as a node-affinity problem and show that their proposed model outperforms alternative dynamic graph-learning approaches. However, a simple exponential-moving-average benchmark achieves performance relatively close to that of the proposed model, highlighting the strong persistence of institutional portfolios and the importance of comparison with simple temporal benchmarks. Their study is particularly relevant to the present thesis because it demonstrates that Form 13F holdings can be represented and modeled directly as a dynamic graph.

A related application by \textcite{tagndynamic2026} constructs a multi-relational financial network combining price correlations, supply-chain relationships, and institutional co-holdings. The proposed temporal graph-attention framework integrates graph attention with recurrent learning to predict extreme stock-price volatility. Institutional co-holdings therefore enter the model as one of several potential channels of risk transmission, and the results indicate that the combined network improves extreme-volatility prediction relative to the study's non-graph benchmarks.

Together, these studies demonstrate the potential of graph learning in applications involving institutional ownership and financial risk, but they do not resolve the specific question examined in this thesis. \textcite{izadifar2026institutional} isolate the institutional ownership graph but predict future portfolio allocations rather than risk outcomes. Conversely, \textcite{tagndynamic2026} predict extreme volatility using several network layers simultaneously, preventing the contribution of institutional co-holdings from being identified separately. It therefore remains unclear whether institutional ownership topology, considered as a distinct source of information, improves the out-of-sample prediction of stock-level downside risk beyond market characteristics and conventional ownership measures, and whether learning directly from the graph provides an advantage over tabular models using engineered network features.
